\documentclass[11pt]{article}
% Shared print layout. The web's provisional design and MathJax are independent.
\usepackage[utf8]{inputenc}
\usepackage[T1]{fontenc}
\usepackage[spanish,es-nodecimaldot]{babel}
\usepackage{lmodern}
\usepackage{amsmath,amssymb}
\usepackage{graphicx}
\usepackage{geometry}
\geometry{a4paper,left=20mm,right=20mm,top=25mm,bottom=22mm,headheight=16pt,headsep=8mm}
\usepackage{fancyhdr,lastpage}
\usepackage{enumitem}
\usepackage{xcolor}
\usepackage{pgfplots}
\pgfplotsset{compat=1.18}
\usepgfplotslibrary{fillbetween}
\usepackage[hidelinks,pdfencoding=auto]{hyperref}
\hypersetup{pdflang=es-ES,pdfauthor={},pdfsubject={Matemáticas Aplicadas a las Ciencias Sociales II}}
\pagestyle{fancy}
\fancyhf{}
\fancyhead[L]{\small Matemáticas Aplicadas a las Ciencias Sociales II}
\fancyhead[R]{\small 2.\textsuperscript{o} Bachillerato}
\fancyfoot[L]{\footnotesize Subida de nota de curso | Edición revisada 1.0}
\fancyfoot[R]{\footnotesize Página \thepage\ de \pageref{LastPage}}
\renewcommand{\headrulewidth}{0.4pt}
\setlength{\parindent}{0pt}
\setlength{\parskip}{0.5em}
\setlist{itemsep=0.3em,topsep=0.3em}
\newcommand{\apartado}[1]{\subsection*{#1}}
\pgfplotsset{
  revision/.style={
    width=0.94\linewidth,height=7.5cm,
    axis lines=middle,grid=major,grid style={gray!18},
    tick label style={font=\small},label style={font=\small},
    xlabel={$x$},ylabel={$y$},samples=160,
    legend style={at={(0.5,-0.16)},anchor=north,draw=none,font=\small},
    clip=true
  }
}

\hypersetup{pdftitle={Subida de nota de curso - Examen revisado}}
\begin{document}
{\Large\bfseries Prueba de subida de nota de curso}

IES ALANDALUS - Departamento de Matemáticas.

Apellidos y nombre: \hrulefill

Fecha: \hrulefill

Duración: \textbf{1 hora y 30 minutos}. Cuatro ejercicios de 2,5 puntos: \textbf{10 puntos}. Responde en el papel de examen, no en los folios de enunciados. Si se responde a dos ejercicios de un bloque, solo se corrige el que aparezca físicamente primero; se conserva la regla original aunque aquí cada bloque contiene un ejercicio.

Se permiten calculadoras no programables, no gráficas y sin capacidad para almacenar o transmitir datos. Justifica todos los procesos. Cada ejercicio incluye 0,25 puntos de expresión correcta de procesos y métodos, sin sumarlos a su máximo.
\section*{Bloque A - Ejercicio 1}
 \textbf{(2,5 puntos)} Sea la función \( f(x) \) definida por:

\[
f(x) =
\begin{cases} 
x^2 - 4x + a, & \text{si } x < 2 \\[8pt]
\displaystyle \frac{1}{x - 1}, & \text{si } x \geq 2
\end{cases}
\]

\begin{enumerate}
    \item[\textbf{a)}] \textbf{(1.3 puntos)} Calcule el valor de \( a \) para que la función sea continua en \( x = 2 \). Para ese valor de \( a \) obtenido, \textquestiondown{}es derivable la función en \( x = 2 \)?
    \item[\textbf{b)}] \textbf{(1.2 puntos)} Para \( a = 4 \), estudie la monotonía y calcule las ecuaciones de las asíntotas, si existen.
\end{enumerate}
\clearpage
\section*{Bloque B - Ejercicio 2}
 \textbf{(2,5 puntos)} Una confitería elabora dos tipos de tartas, unas de chocolate y otras de merengue y chocolate. Para ello dispone de 100 kg de bizcocho, 80 kg de crema de chocolate y 46 kg de merengue. Para elaborar una tarta de chocolate, se requieren 1 kg de bizcocho y 2 kg de crema de chocolate y para la tarta de chocolate y merengue se requieren 2 kg de bizcocho, 1 kg de crema de chocolate y 1 kg de merengue. Por cada tarta de chocolate se obtiene un beneficio de 10 euros y de 12 euros por cada una de merengue y chocolate. Suponiendo que se vende todo lo que se elabora, \textquestiondown{}cuántas tartas de cada tipo debe preparar para obtener un beneficio máximo? \textquestiondown{}Cuál es dicho beneficio?


\clearpage
\section*{Bloque C - Ejercicio 3}
 \textbf{(2,5 puntos)} Sean \( A \) y \( B \) dos sucesos de un mismo experimento aleatorio.

\begin{enumerate}
    \item[\textbf{a)}] \textbf{(0.5 puntos)} Si \( P(A) \neq 0 \) y \( P(B) \neq 0 \), \textquestiondown{}pueden ser los sucesos \( A \) y \( B \) independientes e incompatibles a la vez? Justifique la respuesta.

    \item[\textbf{b)}] \textbf{(2 puntos)} Sabiendo que \( P(A) = 0{,}3 \), \( P(B) = 0{,}5 \) y \( P(A/B) = 0{,}2 \), calcule las siguientes probabilidades:
    \[
    P(A \cap B) \qquad P(A \cup B) \qquad P(A^c \cup B^c) \qquad P(A - B)
    \]
\end{enumerate}

\clearpage
\section*{Bloque D - Ejercicio 4}
 \textbf{(2,5 puntos)} El tiempo de desfase, en minutos, entre la hora de paso programada de un autobús por cierta parada y la hora real a la que pasa, sigue una distribución Normal de media desconocida y varianza 4. Se observa el paso del autobús por la parada en 10 ocasiones elegidas al azar, registrándose los siguientes desfases:

\[
4.7 \quad 2.1 \quad 3.6 \quad 5.4 \quad 0.0 \quad 4.2 \quad 4.0 \quad -0.2 \quad 1.9 \quad 5.2
\]

\begin{itemize}
    \item[a)] \textbf{(1.25 puntos)} Obtenga un intervalo de confianza al 97\% para el desfase medio en la hora de paso del autobús.

    \item[b)] \textbf{(1.25 puntos)} \textquestiondown{}Qué tamaño muestral mínimo sería necesario para estimar el desfase medio con un error inferior a 30 segundos y un nivel de confianza del 95\%? \textquestiondown{}Cómo variaría dicho tamaño muestral si se aumentara el nivel de confianza?
\end{itemize}



\end{document}
